\documentclass[11pt,a4paper]{article}
% Shared preamble of the FEBAUTO Racing manuals.
\usepackage[margin=2.2cm,headheight=16pt]{geometry}
\usepackage[T1]{fontenc}
\usepackage{lmodern}
\usepackage{xcolor}
\usepackage{graphicx}
\usepackage{booktabs}
\usepackage{tabularx}
\usepackage{listings}
\usepackage{fancyhdr}
\usepackage{enumitem}
\usepackage{tikz}
\usetikzlibrary{arrows.meta}
\usepackage[hidelinks,colorlinks=true,linkcolor=febblue,urlcolor=febblue]{hyperref}

\definecolor{febblue}{HTML}{1C4687}
\definecolor{febgold}{HTML}{FDB515}
\definecolor{codebg}{HTML}{F3F4F6}

\lstset{
  basicstyle=\ttfamily\small,
  backgroundcolor=\color{codebg},
  frame=single, framerule=0pt, framesep=4pt,
  xleftmargin=4pt, xrightmargin=4pt,
  breaklines=true, columns=fullflexible, keepspaces=true,
  showstringspaces=false,
}
\newcommand{\file}[1]{\texttt{#1}}
\newcommand{\cmd}[1]{\texttt{#1}}
\newcommand{\rulebox}[1]{\par\noindent\colorbox{codebg}{\parbox{\dimexpr\linewidth-2\fboxsep}{#1}}\par\medskip}

\pagestyle{fancy}
\fancyhf{}
\fancyhead[L]{\color{febblue}\textbf{FEBAUTO Racing}}
\fancyhead[R]{\color{febblue}\nouppercase{\leftmark}}
\fancyfoot[C]{\thepage}
\renewcommand{\headrulewidth}{0.8pt}
\renewcommand{\headrule}{\color{febgold}\hrule width\headwidth height\headrulewidth}
\setcounter{tocdepth}{1}

% \febtitle{subtitle}{part lines}
\newcommand{\febtitle}[2]{%
\begin{titlepage}
  \centering
  \vspace*{2cm}
  \includegraphics[width=5cm]{../../assets/feb_logo_512.png}\\[1.5cm]
  {\color{febblue}\Huge\bfseries FEBAUTO Racing}\\[6pt]
  {\Large #1}\\[2cm]
  \begin{tabular}{rl}#2\end{tabular}\\[3cm]
  Formula Electric at Berkeley\\
  September 2026, platform release v0.1.0\\[6pt]
  \url{https://github.com/FEBAutonomous/feb-racing}\\
  \url{https://febautonomous.github.io/feb-racing/}
\end{titlepage}
\tableofcontents
\newpage}

\setlist{nosep,leftmargin=*}

\begin{document}
\febtitle{Lead's car manual}{%
  \textbf{Part I} & Power, networks, a terminal on the car\\[3pt]
  \textbf{Part II} & Running a driver, the pad, what to watch\\[3pt]
  \textbf{Part III} & When something is wrong, from what actually went wrong\\}

\noindent{\small\color{gray}Everything here was done on the car between 2026-10-18 and 2026-10-24. Member-facing
material is in \file{member\_guide.pdf}; this is the rest.}

\part{Power, networks, a terminal}

\section{Power on, in this order}

\begin{enumerate}
  \item Gamepad on. Hold its power button until the LED is solid. \emph{Before} the stack starts: the joystick
  node looks for the pad once, at start.
  \item Battery in, balance plug into the power board, SW1 on. The Jetson boots off the power board; about a
  minute. If the Jetson is on its USB-C brick instead, SW1 only powers the VESC.
  \item A terminal on the car (next section), then the stack (Part II).
  \item Off: Ctrl-C the stack, SW1 off, battery out. The Jetson does not need a clean shutdown.
\end{enumerate}

Battery: a 3S LiPo, 12.6~V full, charge below 11.1, never below 10.5. The VESC reports the voltage; the
power board runs warm, the battery must not. Store it in the bag, never in the car.

\section{Reaching the car}

The Jetson's login is \texttt{feb}, password \texttt{feb}. In the order it tries networks at boot:

\begin{tabularx}{\linewidth}{lX}
\toprule
Where & How \\
\midrule
home Wi-Fi \texttt{mandem}, or the phone hotspot \texttt{pranavphone} & it joins by itself (the hotspot needs the Personal Hotspot page open on the phone until it connects). Then \cmd{ssh feb@feb-car} from any laptop on the tailnet, anywhere with internet. \\
nothing known around & 45~s after boot it becomes Wi-Fi network \texttt{feb-car}, password \texttt{febautonomous}. Join it, \cmd{ssh feb@10.42.0.1}. No internet for the car in this mode. \\
USB-C cable to a laptop & \cmd{ssh feb@192.168.55.1} \\
\bottomrule
\end{tabularx}

\medskip
Gotcha: if the car boots without a known Wi-Fi it goes into its own hotspot and \emph{stays} there; it will not
join the phone later. Either power-cycle with the hotspot page open, or join \texttt{feb-car} and run
\cmd{sudo nmcli con up pranavphone}. Adding a new Wi-Fi:

\begin{lstlisting}
sudo nmcli con add type wifi ifname wlP1p1s0 con-name NAME ssid NAME wifi-sec.key-mgmt wpa-psk wifi-sec.psk PASSWORD connection.autoconnect yes
\end{lstlisting}

The car's Wi-Fi region is set to US (5~GHz works). Tailscale name \texttt{feb-car}, plain ssh, Tailscale SSH off.

\part{Running a driver}

\section{The command}

\begin{lstlisting}
cd ~/FEB/feb-stack && git pull                      # the team's merged drivers
cd ~/FEB/feb-racing && git pull                     # launcher, tracks, docs
./feb up --car --driver feb_driver --mode-name inspection
\end{lstlisting}

Drivers come from \file{feb-stack} first, then the starter kit, so the starter drivers run before anyone
has submitted. \cmd{-{}-mode-name inspection} caps the supervisor at 0.8~m/s; leave it off for the driver's
trackdrive cap. Wait for three lines in the output: \emph{Connected to VESC}, \emph{Opened joystick},
\emph{scan\_rotate: lidar\_yaw}. Then the pad.

\begin{tabularx}{\linewidth}{lX}
\toprule
Driver & What it does on the car \\
\midrule
\texttt{feb\_driver} & the gap follower; drove the corridor on 2026-10-19 at 0.8~m/s, door nooks included \\
\texttt{feb\_racer} & a mapping lap on the gap-follower law at 1.2~m/s, then a raceline at up to 2.5~m/s in trackdrive. Untried on the car as of 2026-10-24: first run in inspection, on a loop it can close, someone walking beside it. \\
\texttt{feb\_cone\_driver} & cones; needs a cone track. \\
\bottomrule
\end{tabularx}

\section{The pad}

\begin{tabularx}{\linewidth}{lX}
\toprule
Hold & \\
\midrule
LB & MANUAL: left stick forward/back is speed, right stick is steering \\
RB & AUTONOMOUS: the driver drives \\
nothing & IDLE, stopped \\
minus & STOP, latched; LB or RB clears it \\
\bottomrule
\end{tabularx}

\medskip
The 8BitDo Ultimate 2C has two modes with different button numbers. The car is configured for the mode it is in
(usb id \texttt{2dc8:301c}: LB 6, RB 7, minus 10). \textbf{Never press the X, B or A + Home combos}; that flips
the mode and every button moves. If a bumper does nothing or STOP latches by itself, that is what happened:
power the pad off, hold X + Home until it lights, check with \cmd{lsusb | grep 8BitDo} that it reads
\texttt{301c}.

\section{What to watch, in a second terminal}

\begin{lstlisting}
./feb logs -f                                         # everything the stack prints
./feb shell -c "ros2 topic hz /scan"                  # 10 Hz from the lidar
./feb shell -c "ros2 topic echo /feb/state"           # IDLE / MANUAL / AUTONOMOUS / STOP and why
./feb shell -c "ros2 topic echo /drive"               # what the driver asks for
\end{lstlisting}

If \cmd{ros2 topic} refuses with \emph{context is invalid} or \emph{failed to initialize wait set}, that is a
quirk of the ROS tools after many short commands; the stack is fine. \cmd{ros2 daemon stop} and try again,
or read \cmd{./feb logs}.

Recording a run: \cmd{./feb shell -c "ros2 bag record -s mcap -o /home/feb/FEB/bags/NAME /scan /drive
/ackermann\_cmd /feb/state /odom /joy /sensors/core"} in the second terminal, Ctrl-C to end; the folder
lands in \file{\textasciitilde/FEB/bags} on the Jetson, \cmd{scp -r} it to a laptop and open the \file{.mcap}
in Foxglove.

\part{When something is wrong}

\begin{tabularx}{\linewidth}{lX}
\toprule
Symptom & Cause and fix \\
\midrule
no \emph{Opened joystick}, or \cmd{ls /dev/input/js0} fails & the pad is off, or paired after the stack started. Pad on, LED solid, then Ctrl-C and \cmd{feb up} again. The dongle shows \texttt{8BitDo IDLE} in \cmd{lsusb} while the pad is off. \\
STOP latches on its own; LB or RB does nothing & the pad changed mode (see The pad). \\
no \emph{Connected to VESC}, no \file{/dev/sensors/vesc} & SW1 off, or the VESC USB in the wrong place. The VESC needs battery power; the Jetson's brick is not enough. \\
\file{/dev/sensors/rplidar} missing & the lidar's USB; the udev rule maps the CP2102 to it. \\
car drives in the wrong direction / steers wrong & VESC Tool: Invert Motor Direction (the motor), servo 0.3 should steer left. Both set on 2026-10-19. \\
lidar scan is rotated & \texttt{lidar\_yaw} in \file{devkit/feb\_tools/config/car.yaml}, 1.5708 for the S2 turned 90 degrees. Check: a hand in front of the nose reads 0 degrees in \cmd{/scan}. \\
Docker: \emph{Unable to enable DIRECT ACCESS FILTERING, raw table} & the Jetson kernel has no iptables raw module; dockerd runs with \texttt{DOCKER\_INSECURE\_NO\_IPTABLES\_RAW=1} in \file{/etc/systemd/system/docker.service.d/raw.conf}. Already in place; re-create after a reflash. \\
\cmd{feb update} replaced the box with the registry image and something is missing & on the car the registry image may lag the local build. Rebuild natively (below) or \cmd{docker tag} the local image back to \texttt{:gui} and \cmd{./feb reset}. \\
a house-code change must reach the car now & \cmd{docker build -t ghcr.io/febautonomous/feb-devkit:latest devkit/} then \cmd{-{}-build-arg BASE=... -f devkit/Dockerfile.gui -t ...:gui}; ten minutes from cache. Then \cmd{./feb reset}. \\
a car.yaml change must reach the car now & \cmd{docker cp} it into the box at \file{/home/autodrive\_devkit/install/feb\_tools/share/feb\_tools/config/car.yaml}, \cmd{feb stop}, \cmd{feb up}. Rebuild the image when convenient. \\
sudo needed from a Claude session & it cannot; run it yourself over ssh with \cmd{-t}. \\
car not on the phone hotspot at campus & it fell back to its own hotspot at boot (Part I). \\
motor stutters from standstill & the sensorless dead zone below ~1000~ERPM; normal. \\
\bottomrule
\end{tabularx}

\section{Reflashing}

Jetson Linux R36.4.3 from this PC with \file{\textasciitilde/FEB/jetson-flash/prepare\_and\_flash.sh}: power off,
a jumper across J14 pins 9 and 10 while plugging the power in, USB-C to the PC, run the script. The Jetson
shows as USB \texttt{0955:7523 APX} in recovery. The script pre-seeds the user, hostname, udev rules and the
hotspot fallback; Tailscale, Docker's raw-table setting, the xpad driver (compiled from source for the
8BitDo receiver, \file{\textasciitilde/xpad-src}) and the Wi-Fi networks are done by hand afterwards:
\file{docs/build-car.md}, Session 4, has each command.

\end{document}
